\section{向量值哈希引理}\label{sec:hashing}

下面的引理把特征之间的弥散转化为表项之间的弥散。其常数与特征空间
的维数以及被哈希向量的维数无关。在应用中，$q$ 将索引候选标记，
$z$ 将索引门表的表项。所得的哈希向量将收集对一个表项的影响有贡献
的全部傅里叶系数，故其范数平方即为该影响。这就是我们需要向量值
命题而非标量估计的原因。

\begin{lemma}[向量值哈希]\label{lem:hashing}
设 $H$ 是 $\F$ 上的有限维向量空间，对 $q\in H^*$ 与 $u\in H$ 记
$\chi_q(u)=(-1)^{q(u)}$；设 $\mathcal V$ 是实希尔伯特空间，
$s\ge1$ 为整数，且
\[
 c(q,z)\in\mathcal V,
 \qquad q\in H^*,\quad z\in\F^s.
\]
令
\[
 w_q=\sum_{z\in\F^s}\norm{c(q,z)}^2,
 \qquad M=\sum_{q\in H^*}w_q.
\]
独立均匀地选取 $L,J_1,\ldots,J_s\in H$，并定义
\[
 \mathbf J(q)=(q(J_1),\ldots,q(J_s)),
 \qquad
 U(z)=\sum_{q\in H^*}\chi_q(L)c(q,z+\mathbf J(q)).
\]
则
\begin{equation}\label{eq:hashing-total-mass}
 \E\sum_{z\in\F^s}\norm{U(z)}^2=M.
\end{equation}
对每个 $\zeta>0$，存在只依赖 $\zeta$ 的 $\gamma>0$ 与
$s_0\in\mathbb N$，使得对每个 $s\ge s_0$，
\begin{equation}\label{eq:hashing-dispersion}
 \max_{q\in H^*}w_q\le\gamma M
 \quad\Longrightarrow\quad
 \E\max_{z\in\F^s}\norm{U(z)}^2\le\zeta M.
\end{equation}
\end{lemma}

\begin{proof}
下文对 $H$ 与 $\F^s$ 的期望都使用均匀概率测度；定义 $w_q$ 与 $M$
的显示求和不作归一化。在给定 $J_g$ 的条件下，$L$ 的特征的正交性给出
\[
 \E_L\norm{U(z)}^2
 =\sum_q\norm{c(q,z+\mathbf J(q))}^2.
\]
对 $z$ 求和即证 \eqref{eq:hashing-total-mass}。若 $M=0$，则每个向量
$c(q,z)$ 为零，\eqref{eq:hashing-dispersion} 平凡成立。以下设
$M>0$。

令 $h=2^s$，$\rho=\max_q w_q/M$。下面对每个 $C,D>0$ 证明定量界，
再选取其参数。主要估计把最大范数平方用四次幂之和的平方根来控制。
经傅里叶换元后，该四阶矩中的大多数项迫使所有四个因子取同一特征
$q$；其总和由特征质量的弥散控制。两次截断将控制剩余项以及换元给
阵列带来的代价：
\begin{equation}\label{eq:hashing-quantitative}
 \E\max_z\norm{U(z)}^2
 \le 2M\left[
 \sqrt{2\rho+\frac{2C^2}{D^2}+\frac{3D^4}{h}}
 +\left(\frac1{\sqrt C}+\frac1D\right)^2+\frac1h
 \right].
\end{equation}


\paragraph{傅里叶表示。}
对 $p\in\F^s$ 与 $u\in H$，定义
\[
 \alpha_q(p)=\sum_{z\in\F^s}(-1)^{p\cdot z}c(q,z),
 \qquad
 A_p(u)=\sum_{q\in H^*}\chi_q(u)\alpha_q(p),
 \qquad
 S_p=L+\sum_{g=1}^s p_gJ_g.
\]
在 $z$ 上作傅里叶反演得
\begin{equation}\label{eq:hashing-inversion}
 U(z)=\E_p(-1)^{p\cdot z}A_p(S_p).
\end{equation}
记 $m_p=\E_u\norm{A_p(u)}^2$，两个群上的 Parseval 恒等式给出
\begin{equation}\label{eq:hashing-parseval}
 m_p=\sum_q\norm{\alpha_q(p)}^2,
 \qquad \E_p m_p=M,
 \qquad \E_p\norm{\alpha_q(p)}^2=w_q.
\end{equation}
对不同的 $p,p'$，变量 $S_p,S_{p'}$ 独立且均匀于 $H$：系数向量
$(1,p),(1,p')$ 在 $\F$ 上线性无关。同理，任意三个不同的
$p_1,p_2,p_3$ 给出三个独立均匀变量 $S_{p_1},S_{p_2},S_{p_3}$。
事实上，其系数向量的任一零和非空子集必有偶数个元素，因而由两个
相等向量组成。

\paragraph{两次截断。}
我们先丢弃平均范数平方 $m_p$ 异常大的指标 $p$，再在其余指标上丢弃
$A_p(u)$ 的异常大值。第一个截断使第二个截断下每个特征系数的改变
一致地小；第二个截断给出四阶矩指标重合时所需的逐点界。我们将用一
阶矩界与由样本 $S_p$ 两两独立得到的方差界来控制被丢弃部分。定义
\[
 A'_p(u)=
 \begin{cases}
 A_p(u),&m_p\le CM\text{ 且 }\norm{A_p(u)}\le D\sqrt M,\\
 0,&\text{否则},
 \end{cases}
\]
并令
\[
 \alpha'_q(p)=\E_u\chi_q(u)A'_p(u),
 \qquad w'_q=\E_p\norm{\alpha'_q(p)}^2,
 \qquad U'(z)=\E_p(-1)^{p\cdot z}A'_p(S_p).
\]
逐点截断不增加总范数平方，故 Parseval 给出 $\sum_qw'_q\le M$。若
$m_p\le CM$，则
\[
 \norm{\alpha_q(p)-\alpha'_q(p)}
 \le\E_u\norm{A_p(u)}
       \one_{\{\norm{A_p(u)}>D\sqrt M\}}
 \le\frac{m_p}{D\sqrt M}
 \le\frac{C\sqrt M}{D}.
\]
若 $m_p>CM$，则 $\alpha'_q(p)=0$。因此
\begin{equation}\label{eq:hashing-truncated-mass}
 \sum_qw'_q\le M,
 \qquad
 w'_q\le 2w_q+\frac{2C^2M}{D^2}.
\end{equation}

令
\[
 X_p=\norm{(A_p-A'_p)(S_p)},\qquad X=\E_pX_p.
\]
Equation~\eqref{eq:hashing-inversion} 与三角不等式蕴含
$\max_z\norm{U(z)-U'(z)}\le X$。$m_p>CM$ 的指标对 $\E X$ 的贡献
至多为
\[
 \E_p\one_{\{m_p>CM\}}\sqrt{m_p}
 \le\frac{\E_p m_p}{\sqrt{CM}}=\sqrt{M/C}.
\]
其余指标的贡献至多 $\E_p m_p/(D\sqrt M)=\sqrt M/D$。此外，由
$S_p$ 两两独立知 $X_p$ 两两独立。由于 $A_p-A'_p$ 在每点要么为
$A_p$ 要么为零，
\[
 \operatorname{Var}(X)
 =h^{-2}\sum_p\operatorname{Var}(X_p)
 \le h^{-2}\sum_p\E_u\norm{A_p(u)-A'_p(u)}^2
 \le\frac Mh.
\]
我们已证明
\begin{equation}\label{eq:hashing-truncation-error}
 \E\max_z\norm{U(z)-U'(z)}^2
 \le M\left[\left(\frac1{\sqrt C}+\frac1D\right)^2+\frac1h\right].
\end{equation}

\paragraph{四阶矩。}
为简洁记 $B_p=A'_p(S_p)$。展开范数的四次幂并对 $z$ 中的特征求和，
得精确恒等式
\begin{align}
 \E\sum_z\norm{U'(z)}^4
 &=h^{-3}\sum_{p_1,p_2,p_3}
   \E\bigl[\ip{B_{p_1}}{B_{p_2}}
                  \ip{B_{p_3}}{B_{p_1+p_2+p_3}}\bigr]\notag\\
 &=\E_{p_1,p_2,p_3}\E
   \bigl[\ip{B_{p_1}}{B_{p_2}}\ip{B_{p_3}}{B_{p_4}}\bigr],
 \qquad p_4=p_1+p_2+p_3.
 \label{eq:hashing-fourth-expansion}
\end{align}
第二行中 $p_1,p_2,p_3$ 是独立均匀指标。第一行的因子 $h^{-3}$ 来自
傅里叶反演中的四个因子 $h^{-1}$ 以及 $h$ 个特征对 $z$ 的求和。

若 $p_1,p_2,p_3,p_4$ 中任意两个重合，则 $p_1,p_2,p_3$ 中某两个
重合。该事件的概率至多 $3/h$。由 $\norm{B_p}\le D\sqrt M$，其对
\eqref{eq:hashing-fourth-expansion} 的贡献绝对值至多 $3D^4M^2/h$。

现在固定四个不同指标并设 $p_4=p_1+p_2+p_3$。前三个 $S_{p_i}$ 是
$H$ 中独立均匀元素，而 $S_{p_4}=S_{p_1}+S_{p_2}+S_{p_3}$。把四个
$A'_{p_i}$ 按特征展开并对这三个独立元素取平均，得
\begin{equation}\label{eq:hashing-common-character}
 \E\bigl[\ip{B_{p_1}}{B_{p_2}}\ip{B_{p_3}}{B_{p_4}}\bigr]
 =\sum_q
   \ip{\alpha'_q(p_1)}{\alpha'_q(p_2)}
   \ip{\alpha'_q(p_3)}{\alpha'_q(p_4)}.
\end{equation}
具体地，由 $q_1,q_2,q_3,q_4$ 索引的特征产生因子
\[
 \prod_{i=1}^3\chi_{q_i+q_4}(S_{p_i}),
\]
除非 $q_1=q_2=q_3=q_4$，其期望为零。

\eqref{eq:hashing-common-character} 的右端不必非负。为得上界，把
每个被加项换成非负的 $\prod_{i=1}^4\norm{\alpha'_q(p_i)}$，再去掉
四指标互异这一限制。对每个固定的 $q$，Cauchy--Schwarz 给出
\begin{align*}
 \E_{p_1,p_2,p_3}\prod_{i=1}^4\norm{\alpha'_q(p_i)}
 &\le
 \left(\E\norm{\alpha'_q(p_1)}^2
             \norm{\alpha'_q(p_2)}^2\right)^{1/2}
 \left(\E\norm{\alpha'_q(p_3)}^2
             \norm{\alpha'_q(p_4)}^2\right)^{1/2}\\
 &=(w'_q)^2.
\end{align*}
最后一个等式用了每对 $(p_1,p_2)$ 与 $(p_3,p_4)$ 的无条件独立均匀
律；它不断言在以互异为条件之后的独立性。于是非退化四阶矩被
$\sum_q(w'_q)^2$ 控制：特征之间的弥散就此控制表项之间的集中。把这些
估计与 \eqref{eq:hashing-truncated-mass} 结合得
\begin{align*}
 \E\sum_z\norm{U'(z)}^4
 &\le \sum_q(w'_q)^2+\frac{3D^4M^2}{h}\\
 &\le M^2\left(2\rho+\frac{2C^2}{D^2}+\frac{3D^4}{h}\right).
\end{align*}
再用一次 Cauchy--Schwarz，
\begin{equation}\label{eq:hashing-truncated-maximum}
 \E\max_z\norm{U'(z)}^2
 \le\left(\E\sum_z\norm{U'(z)}^4\right)^{1/2}
 \le M\sqrt{2\rho+\frac{2C^2}{D^2}+\frac{3D^4}{h}}.
\end{equation}
不等式 $(a+b)^2\le2a^2+2b^2$ 连同
\eqref{eq:hashing-truncation-error} 与
\eqref{eq:hashing-truncated-maximum} 证明
\eqref{eq:hashing-quantitative}。

最后令 $\varepsilon=\min\{\zeta,1\}$，并选
\[
 C=\frac{64}{\varepsilon},\qquad
 D=\frac{8C}{\varepsilon},\qquad
 \gamma=\frac{\varepsilon^2}{128},\qquad
 2^{s_0}\ge
 \max\left\{\frac{192D^4}{\varepsilon^2},\frac{16}{\varepsilon}\right\}.
\]
若 $s\ge s_0$ 且 $\rho\le\gamma$，则
\eqref{eq:hashing-quantitative} 中的平方根项至多 $\varepsilon/4$。
又 $D\ge\sqrt C$，故其括号中另两项各至多 $\varepsilon/16$。所得界
为 $3\varepsilon M/4\le\zeta M$，以只依赖 $\zeta$ 的常数建立
\eqref{eq:hashing-dispersion}。
\end{proof}
